\begin{table}[H]
\caption{\label{tab:Construction-of-task}Construction of task contents measures}

\begin{tabular}{>{\raggedright}m{0.25\textwidth}>{\raggedright}m{0.68\textwidth}}
\toprule 
Task Content Measure (\textbf{\textit{T}}) & Task Items (\textbf{\textit{J}})\tabularnewline
\midrule Non-routine cognitive analytical & Analyzing data/information (4.A.2.a.4)\tablefootnote{Codes in brackets corresponds to O{*}NET Work Activities, Work Context,
Abilities and Skills Importance scales. {[}\href{https://www.onetcenter.org/database.html\#all-files}{https://www.onetcenter.org/database.html\#all-files}{]}}

Thinking creatively (4.A.2.b.2)

Interpreting information for others (4.A.4.a.1)\tabularnewline
\midrule Non-routine cognitive interpersonal & Establishing and maintaining personal relationships (4.A.4.a.4)

Guiding, directing and motivating subordinates (4.A.4.b.4)

Coaching/developing others (4.A.4.b.5)\tabularnewline
\midrule Routine cognitive & The importance of repeating the same tasks (4.C.3.b.7)

The importance of being exact or accurate (4.C.3.b.4)

Structured vs. unstructured work (4.C.3.b.8)\tabularnewline
\midrule Routine manual & Pace determined by the speed of equipment (4.C.3.d.3)

Controlling machines and processes (4.A.3.a.3)

Spending time making repetitive motions (4.C.2.d.1.i)\tabularnewline
\midrule Non-routine manual physical & Operating vehicles, mechanized devices, or equipment (4.A.3.a.4)

Spending time using hands to handle, control or feel objects, tools
or controls (4.C.2.d.1.g)

Manual dexterity (1.A.2.a.2)

Spatial orientation (1.A.1.f.1)\tabularnewline
\bottomrule
\end{tabular} \medskip

\textbf{Notes:} Based on the elaborations of \citet{RePEc:eko:ekoeko:47_23},
who build the structure on the work of \citet{Acemoglu2011}.
\end{table}
